\section{Dataset Construction and Handling}
\label{sec:dataset}

ApexID combines publicly available wildlife individual-identification datasets
that differ in taxonomic coverage, annotation conventions, image quality, and
metadata. We standardize these sources through four stages: data collection and
label harmonization, near-duplicate removal and animal-focused preprocessing,
manual curation, and observation-aware train--validation partitioning. External
test datasets remain fully excluded from model development.

\subsection{Dataset Sources and Scope}
\label{sec:dataset:sources}

We retain images with a reliable individual identity and map each sample to a
common representation containing its source dataset, species or coarser animal
category, identity, and available observation metadata. Identity labels are
namespaced by source to prevent collisions between independently collected
datasets. Table~\ref{tab:apexid_train_datasets} consolidates the composition and
partition statistics of the corpus after all processing; the original releases
remain the authoritative sources for collection and annotation protocols.

\begin{table*}[t]
\centering
\caption{Datasets in the ApexID corpus after preprocessing and manual
curation. ``Train IDs'' counts identities represented in the training
partition; image and observation totals include both training and validation.}
\label{tab:apexid_train_datasets}
\resizebox{\textwidth}{!}{%
\begin{tabular}{l l r r r r r}
\toprule
\textbf{Dataset} & \textbf{Species / Taxon} & \textbf{Train IDs} &
\textbf{\# Images} & \textbf{\# Observations} & \textbf{Train} &
\textbf{Validation} \\
\midrule
AmvrakikosTurtles~\cite{adam2025exploiting} & Sea turtle & 41 & 141 & 141 & 59 & 82 \\
BalearicLizard~\cite{alcaraz2026long} & Lizard & 656 & 4,359 & 4,359 & 2,859 & 1,500 \\
CTai~\cite{freytag2016chimpanzee} & Chimpanzee & 57 & 3,609 & 3,609 & 2,443 & 1,166 \\
CzechLynx~\cite{picek2026czechlynx} & Lynx & 279 & 37,829 & 37,829 & 28,792 & 9,037 \\
GiraffeZebraID~\cite{parham2017animal} & Giraffe; zebra & 876 & 4,333 & 4,333 & 3,569 & 764 \\
Giraffes~\cite{miele2021revisiting} & Giraffe & 136 & 892 & 892 & 651 & 241 \\
HappyWhale~\cite{cheeseman2026happywhale} & Whales; dolphins & 5,538 & 41,407 & 41,407 & 26,304 & 15,103 \\
HyenaID2022~\cite{botswana2022hyena} & Hyena & 211 & 2,422 & 2,422 & 1,599 & 823 \\
IPanda50~\cite{wang2021giant} & Panda & 42 & 5,167 & 5,167 & 3,619 & 1,548 \\
Jaguar\_data~\cite{rueda2025jaguars} & Jaguar & 27 & 1,178 & 1,178 & 795 & 383 \\
LeopardID2022~\cite{botswana2022leopard} & Leopard & 251 & 4,953 & 4,953 & 3,494 & 1,459 \\
NDD20~\cite{trotter2020ndd20} & Dolphin & 68 & 2,122 & 2,122 & 1,488 & 634 \\
NyalaData~\cite{vcermak2024wildlifedatasets} & Nyala & 150 & 1,442 & 1,442 & 951 & 491 \\
ReunionTurtles~\cite{adam2025exploiting} & Sea turtle & 69 & 237 & 237 & 104 & 133 \\
SalamanderID2025~\cite{animalclef2025} & Salamander & 242 & 1,078 & 1,078 & 675 & 403 \\
SealID~\cite{nepovinnykh2022sealID} & Seal & 48 & 1,668 & 1,668 & 1,101 & 567 \\
WhaleSharkID~\cite{holmberg2009estimating} & Whale shark & 388 & 5,958 & 5,958 & 3,855 & 2,103 \\
ZakynthosTurtles~\cite{adam2025exploiting} & Sea turtle & 33 & 114 & 114 & 47 & 67 \\
HumpbackWhaleID~\cite{vcermak2024wildlifedatasets} & Humpback whale & 2,565 & 13,624 & 13,624 & 8,516 & 5,108 \\
\midrule
\textbf{Total} & \textbf{Multi-species} & \textbf{11,677} &
\textbf{132,533} & \textbf{132,533} & \textbf{90,921} &
\textbf{41,612} \\
\bottomrule
\end{tabular}%
}
\end{table*}

The training partition contains 11,677 identities. The validation partition
contains 12,439 identities, comprising 10,768 identities also represented in
training and 1,671 validation-only identities. Observation counts are computed
from the observation-group identifiers rather than directly from image rows;
in the final training and validation partitions, each retained image has a
distinct observation identifier.

\subsection{Preprocessing and Curation}
\label{sec:dataset:curation}

The public releases are first converted to a common directory and annotation
format. Because source file structures and metadata conventions differ, this
step is implemented separately for each dataset before the shared pipeline is
applied.

\paragraph{Near-duplicate removal.}
\label{sec:dataset:duplicates}
Repeated images can encourage memorization and inflate evaluation when
correlated samples cross partitions. We therefore compute MegaDescriptor
embeddings~\cite{vcermak2024wildlifedatasets} and compare images only within
the same identity. If their cosine similarity exceeds the dataset-specific
threshold in Table~\ref{tab:duplicate_thresholds}, one image is removed. We do
not compare different identities because legitimately similar individuals
must remain represented. A threshold of $2.0$ disables filtering for datasets
where it was not applied, since cosine similarity cannot exceed $1$.

\begin{table}[t]
\centering
\caption{Cosine-similarity thresholds for within-identity near-duplicate
removal. A value of $2.0$ denotes disabled filtering.}
\label{tab:duplicate_thresholds}
\resizebox{\linewidth}{!}{%
\begin{tabular}{l r l r}
\toprule
\textbf{Dataset} & \textbf{Threshold} &
\textbf{Dataset} & \textbf{Threshold} \\
\midrule
AmvrakikosTurtles & 0.87 & BalearicLizard & 0.91 \\
CTai & 2.00 & CzechLynx & 0.93 \\
GiraffeZebraID & 0.89 & Giraffes & 0.87 \\
HappyWhale & 2.00 & HyenaID2022 & 0.90 \\
IPanda50 & 0.97 & Jaguar\_data & 2.00 \\
LeopardID2022 & 0.95 & NDD20 & 0.87 \\
NyalaData & 0.96 & ReunionTurtles & 0.95 \\
SalamanderID2025 & 2.00 & SealID & 0.92 \\
WhaleSharkID & 0.93 & ZakynthosTurtles & 0.90 \\
HumpbackWhaleID & 2.00 & -- & -- \\
\bottomrule
\end{tabular}%
}
\end{table}

\paragraph{Animal-focused segmentation.}
\label{sec:dataset:segmentation}
After duplicate removal, we use the pretrained Segment Anything Model~3
(SAM~3)~\cite{carion2025sam3} with a species-level text prompt, such as
``photo of a giraffe.'' The
highest-confidence valid animal mask is applied to the RGB image, and pixels
outside the mask are replaced with black. This common preprocessing reduces
access to location and background cues while retaining the animal appearance
used for re-identification.

\paragraph{Manual curation.}
The processed corpus is then reviewed manually using \textit{V-Rush}, our
in-house data-curation software~\cite{elbassat_vrush_2026}. Reviewers inspect the images for invalid or
poor-quality samples, incomplete segmentations, and dataset-specific inconsistencies. Samples whose identity cannot be resolved reliably are excluded rather than relabeled by inference. The counts in
Table~\ref{tab:apexid_train_datasets} refer to the data retained after this
manual review.

\subsection{Observation-Aware Partitioning}
\label{sec:dataset:partitioning}

An \emph{identity} denotes one annotated animal, whereas an
\emph{observation} groups images from the same encounter, capture event, or
acquisition sequence. Images within an observation can share pose,
illumination, background, and camera conditions and therefore should not be
treated as independent evidence.

Following the MegaDescriptor splitting strategy~\cite{vcermak2024wildlifedatasets},
we use timestamps to construct temporally separated training and validation
observations whenever they are available. Otherwise, DINO image embeddings are
used to group visually similar samples, and each group is assigned wholly to
one partition. The same identity may occur in both partitions only through
independent observations; images in the same timestamp-derived observation or
DINO similarity group remain together.

\subsection{External Test Datasets}
\label{sec:dataset:test}

ELPephants, GCN-ID, BristolGorillas2020, and a private jaguar dataset are
reserved for cross-dataset evaluation and are never used for optimization,
checkpoint selection, or hyperparameter tuning.
Table~\ref{tab:apexid_test_datasets} reports the versions used after
evaluation-specific preprocessing.

\begin{table*}[t]
\centering
\caption{Completely held-out datasets used for external evaluation. Counts
report the samples assigned to gallery, known-query, or novel-query roles in
each of the five evaluation seeds, excluding processed samples that do not
satisfy the role-assignment requirements.}
\label{tab:apexid_test_datasets}
\begin{tabular}{l l r r r}
\toprule
\textbf{Dataset} & \textbf{Species} & \textbf{\# Identities} &
\textbf{\# Images} & \textbf{\# Observations} \\
\midrule
ELPephants~\cite{korschens2019elpephants} & African forest elephant & 272 & 2,049 & 2,028 \\
GCN-ID~\cite{rosa2025gcnid} & Great crested newt & 207 & 1,232 & 1,152 \\
BristolGorillas2020~\cite{brookes2021bristolgorillas2020,brookes2021gorillas} & Western lowland gorilla & 7 & 5,324 & 4,215 \\
Private jaguar test set & Jaguar & 51 & 1,029 & 672 \\
\midrule
\textbf{Total} & -- & \textbf{537} & \textbf{9,634} & \textbf{8,067} \\
\bottomrule
\end{tabular}
\end{table*}

\paragraph{ELPephants.}
ELPephants~\cite{korschens2019elpephants} is used as a single held-out test set,
and the ApexID training corpus contains no elephant images. Table~\ref{tab:elpephants_curation}
summarizes its curation. We remove three unusable images, eight images forming
four duplicate pairs with contradictory identity labels, and one image from
each of 16 consistently labeled duplicate pairs. The final set contains 2,051
images of 274 identities. SAM~3 produces 1,274 masks automatically, while the
remaining 777 masks are refined before evaluation. Of this processed pool,
2,049 images from 272 identities and 2,028 observations satisfy the
evaluation-role requirements used in Table~\ref{tab:apexid_test_datasets}.

Images are segmented using pretrained SAM~3 with the text prompt
``elephant'' and detection and mask thresholds of $0.30$. All detected elephant
instances are merged into a single foreground mask, and pixels outside the
merged mask are replaced with black.

\begin{table}[t]
\centering
\caption{Curation of the ELPephants external test set.}
\label{tab:elpephants_curation}
\resizebox{\linewidth}{!}{%
\begin{tabular}{l r r}
\toprule
\textbf{Stage} & \textbf{Images} & \textbf{Identities} \\
\midrule
ELPephants (original) & 2,078 & 276 \\
$-$ Unusable (ambiguous subject or ear only) & $-3$ & 0 \\
$-$ Duplicate pairs with contradictory labels & $-8$ & $-2$ \\
$-$ Consistent duplicate copies & $-16$ & 0 \\
\midrule
\textbf{Final test set} & \textbf{2,051} & \textbf{274} \\
\quad Automatic SAM~3 masks & 1,274 & -- \\
\quad Refined masks & 777 & -- \\
\bottomrule
\end{tabular}%
}
\end{table}

\paragraph{GCN-ID.}
GCN-ID contains 1,232 images of 207 Great Crested Newts. Each image is
converted to RGB and processed with SAM~3 using the prompt ``newt or
salamander,'' a detection threshold of $0.5$, and a mask threshold of $0.5$.
When multiple masks are returned, we retain the highest-confidence mask. Five
images without a valid prediction are segmented manually rather than removed.

\paragraph{Private jaguar test set.}
This private evaluation dataset is distinct from the
Jaguar\_data source used in the ApexID training corpus. It remains completely
held out from model development and contains 1,040 processed images of 54
individuals, grouped into 675 observations. After applying the evaluation-role
requirements, 1,029 images from 51 individuals and 672 observations are
assigned to the gallery and query roles in each seed.

\paragraph{BristolGorillas2020.}
\label{sec:dataset:bristol}
BristolGorillas2020~\cite{brookes2021bristolgorillas2020,brookes2021gorillas}
contains images of seven western lowland gorillas and is used exclusively for
external evaluation. We preserve its original train--test organization and
identity-folder structure throughout preprocessing.

We apply annotation-guided foreground segmentation using pretrained SAM~3
without additional fine-tuning~\cite{carion2025sam3}. Each image is converted to RGB,
and its YOLO-format face annotation is parsed. SAM~3 is queried sequentially
with the prompts ``gorilla,'' ``western lowland gorilla,'' ``gorilla body,''
and ``ape.'' Candidate masks are filtered using SAM confidence, foreground
area, correspondence with the annotated face location, face coverage, and
containment of the face box. Duplicate candidates are removed based on mask
overlap. If no candidate passes these criteria, a fallback stage prompts SAM~3
with expanded face regions. The selected foreground is retained against a
black background with $10\%$ spatial padding.

Of 5,409 source images, 5,324 remain after excluding invalid annotations and
images without an accepted mask, yielding a retention rate of $98.43\%$. The
retained data comprise 4,238 training and 1,086 test images across all seven
identities. These original partitions are preserved, and the segmented images
are evaluated without further training or identity-specific adaptation.

\subsection{Responsible Data Handling}
\label{sec:dataset:responsible}

All sources remain governed by their original licenses, and ApexID does not
claim ownership of their images. Where redistribution is prohibited, we
release identifiers, preprocessing code, and acquisition instructions instead
of the data. Dataset-level attribution is retained throughout the unified
corpus.

The collection is subject to geographic, taxonomic, and acquisition biases:
species and locations are represented unevenly, and some identities are tied
to particular cameras or environments. We therefore report source-level
statistics and evaluate on independent datasets rather than interpreting an
aggregate result as uniform performance across wildlife populations.
